% GENERATED FILE. Edit canonical lessons, then run npm run generate:books.
% canonical-source-hash: fnv1a64:6ca6fba46695daf4
% canonical-lessons: RU-C300-navodneniye, RU-C300-zemletryaseniye, RU-C300-smerch, RU-C300-zasukha, RU-C300-zatoplyat

\chapter{Natural disasters}
\label{ch:ru-natural-disasters}
\hlchaptermodality{300}

\begin{chapteropening}
\textbf{By the end of this chapter:} \emph{I can talk about a flood, an earthquake, a tornado and a drought, and say that water flooded a place.}

Complete the last lesson of chapter 300: I can talk about a flood, an earthquake, a tornado and a drought, and say that water flooded a place.
\end{chapteropening}

\section[\texorpdfstring{\ru{наводнение} (navodnéniye)}{наводнение (navodnéniye)}]{\ru{наводнение} (navodnéniye) — a flood}

\label{lesson:RU-C300-navodneniye}



\begin{quote}
Before the new one: say the Russian for the person on duty, then the Russian for to call.
\end{quote}

\subsection*{You'll want to know: \ru{наводнение}}
\textbf{\ru{наводнение}} (\emph{navodnéniye}) — \textquotedblleft{}a flood\textquotedblright{}.

\begin{quote}
\ru{После} \ru{дождей} \ru{началось} \ru{наводнение}.
\end{quote}

\begin{quote}
\emph{Pósle dozhdéy nachalós navodnéniye.}
\end{quote}

\begin{quote}
After the rains a flood began.
\end{quote}

Neuter.

\begin{cousinweb}
\emph{\ru{наводнение}} is built on \emph{\ru{наводнить}} \textquotedblleft{}to flood\textquotedblright{}, from \emph{\ru{вода}} \textquotedblleft{}water\textquotedblright{}, inherited from Proto-Slavic \emph{voda} \textquotedblleft{}water\textquotedblright{} (EtymDB); English \emph{water} is a cousin.
\end{cousinweb}

\begin{grammarlens}[title={What you've built}]
One more word for telling, explaining and describing at B1.
\end{grammarlens}

\subsection*{Your turn}
\begin{itemize}
\raggedright
  \item \emph{Say it:} \emph{navodnéniye}
  \item \emph{Say it:} \emph{navodnéniye}, once more
  \item \emph{Say it:} a sentence of your own with \emph{navodnéniye}
  \item \emph{Recall:} say the Russian for the person on duty, then the Russian for to call, then say \emph{navodnéniye} again
\end{itemize}

\begin{tcolorbox}[breakable,colback=teal!4,colframe=teal!35!black,title={Before you move on}]
What does \ru{наводнение} mean? (\textquotedblleft{}A flood\textquotedblright{}.) Say it once more.
\end{tcolorbox}
\section[\texorpdfstring{\ru{землетрясение} (zemletryaséniye)}{землетрясение (zemletryaséniye)}]{\ru{землетрясение} (zemletryaséniye) — an earthquake}

\label{lesson:RU-C300-zemletryaseniye}



\begin{quote}
Before the new one: say the Russian for to call, then the Russian for a flood.
\end{quote}

\subsection*{You'll want to know: \ru{землетрясение}}
\textbf{\ru{землетрясение}} (\emph{zemletryaséniye}) — \textquotedblleft{}an earthquake\textquotedblright{}.

\begin{quote}
\ru{Землетрясение} \ru{было} \ru{ночью}.
\end{quote}

\begin{quote}
\emph{Zemletryaséniye býlo nóchyu.}
\end{quote}

\begin{quote}
The earthquake struck at night.
\end{quote}

Neuter.

\begin{cousinweb}
\emph{\ru{землетрясение}} joins \emph{\ru{земля}} \textquotedblleft{}earth\textquotedblright{}, inherited from Proto-Slavic \emph{zemľa}, PIE \emph{d\textsuperscript{h}éǵ\textsuperscript{h}ōm}, and \emph{\ru{трясти}} \textquotedblleft{}to shake\textquotedblright{}, Proto-Slavic \emph{tręsti} (EtymDB): a shaking of the earth.
\end{cousinweb}

\begin{grammarlens}[title={What you've built}]
One more word for telling, explaining and describing at B1.
\end{grammarlens}

\subsection*{Your turn}
\begin{itemize}
\raggedright
  \item \emph{Say it:} \emph{zemletryaséniye}
  \item \emph{Say it:} \emph{zemletryaséniye}, once more
  \item \emph{Say it:} a sentence of your own with \emph{zemletryaséniye}
  \item \emph{Recall:} say the Russian for to call, then the Russian for a flood, then say \emph{zemletryaséniye} again
\end{itemize}

\begin{tcolorbox}[breakable,colback=teal!4,colframe=teal!35!black,title={Before you move on}]
What does \ru{землетрясение} mean? (\textquotedblleft{}An earthquake\textquotedblright{}.) Say it once more.
\end{tcolorbox}
\section[\texorpdfstring{\ru{смерч} (smerch)}{смерч (smerch)}]{\ru{смерч} (smerch) — a tornado, a whirlwind}

\label{lesson:RU-C300-smerch}



\begin{quote}
Before the new one: say the Russian for a flood, then the Russian for an earthquake.
\end{quote}

\subsection*{You'll want to know: \ru{смерч}}
\textbf{\ru{смерч}} (\emph{smerch}) — \textquotedblleft{}a tornado, a whirlwind\textquotedblright{}.

\begin{quote}
\ru{Смерч} \ru{сломал} \ru{много} \ru{деревьев}.
\end{quote}

\begin{quote}
\emph{Smerch slomál mnógo derévyev.}
\end{quote}

\begin{quote}
The tornado broke a lot of trees.
\end{quote}

Masculine; genitive \emph{\ru{смерча}}.

\begin{cousinweb}
\emph{\ru{смерч}} is an old East Slavic word whose origin is not settled; EtymDB lists no parent.
\end{cousinweb}

\begin{grammarlens}[title={What you've built}]
One more word for telling, explaining and describing at B1.
\end{grammarlens}

\subsection*{Your turn}
\begin{itemize}
\raggedright
  \item \emph{Say it:} \emph{smerch}
  \item \emph{Say it:} \emph{smerch}, once more
  \item \emph{Say it:} a sentence of your own with \emph{smerch}
  \item \emph{Recall:} say the Russian for a flood, then the Russian for an earthquake, then say \emph{smerch} again
\end{itemize}

\begin{tcolorbox}[breakable,colback=teal!4,colframe=teal!35!black,title={Before you move on}]
What does \ru{смерч} mean? (\textquotedblleft{}A tornado, a whirlwind\textquotedblright{}.) Say it once more.
\end{tcolorbox}
\section[\texorpdfstring{\ru{засуха} (zásukha)}{засуха (zásukha)}]{\ru{засуха} (zásukha) — a drought}

\label{lesson:RU-C300-zasukha}



\begin{quote}
Before the new one: say the Russian for an earthquake, then the Russian for a tornado.
\end{quote}

\subsection*{You'll want to know: \ru{засуха}}
\textbf{\ru{засуха}} (\emph{zásukha}) — \textquotedblleft{}a drought\textquotedblright{}.

\begin{quote}
\ru{Летом} \ru{была} \ru{сильная} \ru{засуха}.
\end{quote}

\begin{quote}
\emph{Létom bylá sílnaya zásukha.}
\end{quote}

\begin{quote}
In the summer there was a severe drought.
\end{quote}

Feminine; stressed on \emph{\ru{за}-}.

\begin{cousinweb}
\emph{\ru{засуха}} is built on \emph{\ru{за}-} and \emph{\ru{сухой}} \textquotedblleft{}dry\textquotedblright{}, inherited from Proto-Slavic \emph{suxŭ}, from PIE \emph{h\textsubscript{2}sews-} \textquotedblleft{}dry\textquotedblright{} (EtymDB).
\end{cousinweb}

\begin{grammarlens}[title={What you've built}]
One more word for telling, explaining and describing at B1.
\end{grammarlens}

\subsection*{Your turn}
\begin{itemize}
\raggedright
  \item \emph{Say it:} \emph{zásukha}
  \item \emph{Say it:} \emph{zásukha}, once more
  \item \emph{Say it:} a sentence of your own with \emph{zásukha}
  \item \emph{Recall:} say the Russian for an earthquake, then the Russian for a tornado, then say \emph{zásukha} again
\end{itemize}

\begin{tcolorbox}[breakable,colback=teal!4,colframe=teal!35!black,title={Before you move on}]
What does \ru{засуха} mean? (\textquotedblleft{}A drought\textquotedblright{}.) Say it once more.
\end{tcolorbox}
\section[\texorpdfstring{\ru{затоплять} (zatoplyát)}{затоплять (zatoplyát)}]{\ru{затоплять} (zatoplyát) — to flood, to submerge}

\label{lesson:RU-C300-zatoplyat}



\begin{quote}
Before the new one: say the Russian for a tornado, then the Russian for a drought.
\end{quote}

\subsection*{You'll want to know: \ru{затоплять}}
\textbf{\ru{затоплять}} (\emph{zatoplyát}) — \textquotedblleft{}to flood, to submerge\textquotedblright{}.

\begin{quote}
\ru{Вода} \ru{затопила} \ru{улицы}.
\end{quote}

\begin{quote}
\emph{Vodá zatopíla úlitsy.}
\end{quote}

\begin{quote}
The water flooded the streets.
\end{quote}

Imperfective; perfective \emph{\ru{затопить}} (\emph{zatopít}). Careful: \emph{\ru{затопить} \ru{печь}} means to light a stove.

\begin{cousinweb}
\emph{\ru{затоплять}} is built on \emph{\ru{за}-} and \emph{\ru{топить}} \textquotedblleft{}to sink, to drown\textquotedblright{}, inherited from Proto-Slavic \emph{topiti} (EtymDB).
\end{cousinweb}

\begin{grammarlens}[title={What you've built}]
One more word for telling, explaining and describing at B1.
\end{grammarlens}

\subsection*{Your turn}
\begin{itemize}
\raggedright
  \item \emph{Say it:} \emph{zatoplyát}
  \item \emph{Say it:} \emph{zatoplyát}, once more
  \item \emph{Say it:} a sentence of your own with \emph{zatoplyát}
  \item \emph{Recall:} say the Russian for a tornado, then the Russian for a drought, then say \emph{zatoplyát} again
\end{itemize}

\begin{tcolorbox}[breakable,colback=teal!4,colframe=teal!35!black,title={Before you move on}]
What does \ru{затоплять} mean? (\textquotedblleft{}To flood, to submerge\textquotedblright{}.) Say it once more.
\end{tcolorbox}
